\subsubsection{Comparaison de la fenêtre dodécaphasique avec l'inverse du centre}
\label{sec:contraste-ventana-reciproco}

La fenêtre finie \(\alpha_{12}^{(0)}\) de
\eqref{eq:alpha-ventana-cero}, obtenue avec \(c_{13}=0\), admet en outre
la comparaison historique avec l'inverse du centre imprimé de la
constante inverse. Définissons
\[
 R_{2018}:=137.035999084,\qquad
 \alpha_{\mathrm{rec},2018}:=R_{2018}^{-1}.
\]
Dans cette opération, on évalue l'inverse du rationnel décimal
\(R_{2018}\). Cette référence se distingue du centre tabulé
\(0.0072973525693\) utilisé dans le tableau précédent.

\begin{center}
\begin{tabular}{@{}ll@{}}
\toprule
Expression & Évaluation arithmétique\\
\midrule
\(\alpha_{12}^{(0)}\) &
\(0.007297352569283800997285105472380663\)\\
\(\alpha_{\mathrm{rec},2018}\) &
\(0.007297352569283800997285105472380663567972560701\ldots\)\\
\(\alpha_{12}^{(0)}-\alpha_{\mathrm{rec},2018}\) &
\(\simeq-5.679725607012965\times10^{-37}\)\\
\((\alpha_{12}^{(0)})^{-1}-R_{2018}\) &
\(\simeq+1.066588006664435\times10^{-32}\)\\
\(\bigl((\alpha_{12}^{(0)})^{-1}-R_{2018}\bigr)/R_{2018}\) &
\(\simeq+7.783268730800000\times10^{-35}\)\\
\bottomrule
\end{tabular}
\end{center}

La relation décimale précise est
\[
 \alpha_{12}^{(0)}
 =10^{-36}\left\lfloor10^{36}\alpha_{\mathrm{rec},2018}\right\rfloor.
\]
Par conséquent, la coïncidence historique correspond à la troncature à
trente-six positions décimales de l'inverse ; son arrondi à cette
position se termine par \(664\), tandis que la fenêtre se termine par \(663\).
Les différences du tableau quantifient la comparaison entre les deux
nombres. La section infinie \(\alpha_A\) conserve sa propre normalisation
compatible et sa propre frontière, distinctes de celles de cette fenêtre finie.

L'ajustement exprime la constante inverse sous la forme
\(137.035999084(21)\), avec une incertitude-type
\(u(R_{2018})=2.1\times10^{-8}\)\cite{codata2018}.
Les décimales obtenues en inversant son centre imprimé appartiennent à
l'évaluation arithmétique de ce rationnel ; l'incertitude expérimentale
reste celle indiquée par l'ajustement. La comparaison précédente avec le
centre tabulé de \(\alpha\) maintient sa référence et son résidu
normalisé propres.
